\section{Bounded search, finite traces, and changing parameters}
\label{sec:traces}

Bounded search and trace reflection are the workhorses of the framework.  They cover direct finite
algorithms, proof replay, approximation procedures, combinatorial enumeration, and most ``run the method
for sufficiently many steps'' arguments once a concrete fuel has been supplied.

\subsection{Trace packages}

A trace package has input \(x\), request \(q\), packed state \(s\), initializer \(s_0(x,q)\), transition
\(T(i,s,x,q)\), and fuel \(F(x,q)\).  The state after \(i\) steps is defined by primitive recursion:
\[
 s(0,x,q)=s_0(x,q),
 \qquad
 s(Si,x,q)=T(i,s(i,x,q),x,q).
\]
A verifier need not expose the entire sequence as a list.  The recursive symbol itself is the trace, and
the defining equations are available as \(\PRAAx\) instances.

The invariant \(I(i,s(i),x,q)\) must be quantifier-free after bounded formulas have been Booleanized.  It
may contain many fields:
\begin{itemize}
  \item a well-formedness equation for the packed state;
  \item arithmetic bounds on counters and coefficients;
  \item an equality identifying a cached value with a recomputation;
  \item a Booleanized universal statement about all positions already processed;
  \item an error budget, encoded by integer inequalities after clearing denominators;
  \item a resource rank or demand cap.
\end{itemize}
There is no need to force the mathematical prose into one small equation.  A finite conjunction of
equations and decidable inequalities is still quantifier-free.

\begin{theorem}[Finite trace reflection]
\label{thm:trace-reflection}
Suppose the initializer, transition, fuel, and invariant are fixed primitive-recursive terms, and there are
object proofs of the invariant at zero, preservation at a successor step, and a final implication from the
invariant at fuel to \(Q(x,q)\).  Then iPRA proves \(P(x,q)\to Q(x,q)\).
\end{theorem}
\begin{proof}
Package the recursive trace term into the invariant formula.  Apply quantifier-free induction to obtain the
universal invariant.  Eliminate the universal quantifier at the fuel term and cut with the final implication.
Finally discharge the precondition.  This is Theorem~\ref{thm:loop-compile} specialized to the trace.
\end{proof}

\subsection{Fix precision before induction}

An important discipline in several earlier reports is worth making explicit.  The requested precision
\(q\) is a free parameter of the entire induction; it is not an unboundedly quantified object whose needed
input precision is chosen during the proof.  A demand pass first computes all internal precisions from
\(q\).  The forward trace then runs with those fixed values.  Consequently the induction invariant is a
quantifier-free formula over \(i,x,q\), the current state, and the precomputed demands.

This ordering prevents an invalid move sometimes hidden in informal constructive analysis: at stage \(i\),
choose a new precision depending on an existential fact about all future stages.  Such a choice may be
nonconstructive or not primitive recursive.  Backward demand followed by forward replay turns it into a
finite calculation before induction begins.

\subsection{Changing-parameter recursion}

Many algorithms do not iterate a single map with fixed external parameters.  Newton-style methods update a
center and radius; elimination algorithms replace a polynomial system; local-to-global arguments change a
chart; a recursive proof may consume one resource and reset another.  This causes no logical problem if all
changing data are part of the packed state.

Let the mathematical recurrence be
\[
  a_{i+1}=F(i,a_i,p_i),
  \qquad
  p_{i+1}=G(i,a_i,p_i).
\]
Pack \((a_i,p_i)\) as one state and use transition
\[
  T(i,\langle a,p\rangle)=\langle F(i,a,p),G(i,a,p)\rangle.
\]
The invariant refers to both components at stage \(i\).  Primitive recursion permits the successor step to
use the counter, previous state, and fixed original inputs, which is exactly what is required.

The subtle issue is not changing parameters but \emph{future demand}.  If computing \(a_i\) to precision
\(r_i\) requires \(p_i\) to precision \(s_i\), and \(p_i\) itself depends on later outputs, the demand
relations must be solved first.  Section~\ref{sec:demand} gives both acyclic and finite cyclic closure.

\subsection{Proof replay and certificate checking}

A finite trace can carry more than numerical state.  It can replay a certificate.  Suppose a certificate is
a list of instructions, and each instruction has a local soundness theorem.  The state stores the current
instruction pointer, current mathematical data, and a Boolean ``all checks so far passed.''  The invariant
says:
\begin{enumerate}
  \item the pointer is within the certificate length;
  \item the stored data equal the result of executing the prefix;
  \item if the Boolean is \(1\), every executed instruction satisfies its local verifier;
  \item the mathematical relation advertised by the verified prefix holds.
\end{enumerate}
The step theorem uses the local soundness theorem for the next instruction.  At the final pointer, the
invariant implies the public result.

This pattern covers Gaussian elimination certificates, polynomial identities, factorization checks,
finite homology computations, graph paths, finite category diagrams, and proof-carrying numerical
approximations.  It also clarifies the producer--verifier distinction: the producer may be sophisticated or
external, while replay remains a fixed primitive-recursive verifier with an iPRA soundness proof.

\subsection{Search as a trace}

Least bounded search from Theorem~\ref{thm:least-search} is a trace whose state is
\((\mathit{found},\mathit{witness})\).  Complete enumeration is a trace whose state is a bounded list.
Dynamic programming is a trace whose state is a finite table.  Consequently one generic loop compiler
supports all of them.  Application libraries should prove reusable step lemmas for common state updates
rather than restating whole searches.

\subsection{Nested loops and finite products}

A loop of \(A(x)\) outer iterations, each with \(B(x,i)\) inner iterations, can be treated in either of two
ways.

\paragraph{Hierarchical proof.}
For a fixed outer state and index, compile the inner loop to a local contract.  Use that contract in the
outer step proof.  This produces two uses of quantifier-free induction but keeps invariants modular.

\paragraph{Flattened proof.}
Use state \((i,j,s)\), where \(j\) advances until \(B(x,i)\), then resets while \(i\) advances.  A fuel
bound is
\[
  \sum_{i<A(x)}(B(x,i)+1),
\]
which is a bounded primitive-recursive sum.  The invariant includes the relationship between flattened time
and \((i,j)\).  This may produce a faster executable trace but a longer proof.

Both methods are valid in iPRA.  The hierarchical method is normally preferable for formal development,
because each local loop becomes a theorem node in the staged graph.

\subsection{Bounded recursion versus unbounded termination}

A statement that a transition relation is well founded is not by itself a primitive-recursive runtime
bound.  To execute a total primitive-recursive program, the framework requires explicit fuel.  There are
three acceptable sources:
\begin{enumerate}
  \item a direct primitive-recursive numerical bound;
  \item a fixed-dimensional ranking function whose value bounds the run;
  \item a finite state-space bound, so repetition or a terminal state must occur within the number of
        states.
\end{enumerate}
A classical termination proof using an unspecified ordinal or compactness argument must be converted into
one of these finite resources before it can drive a trace.  Section~\ref{sec:resources} develops the
fixed-dimensional lexicographic case.

\subsection{Uniformity and specialization}

There are two distinct notions of uniformity.

\begin{itemize}
  \item A theorem can be uniform in numerical inputs: one fixed primitive-recursive program works for all
        \(x,q\).
  \item A compiler can be uniform over a family of finite syntax descriptions at the Lean meta-level: for
        each description it generates a different concrete primitive-recursive program and proof.
\end{itemize}
The second does not imply that iPRA contains one evaluator for all descriptions.  For example, a matrix
algorithm may be generated for every fixed dimension \(d\), producing a family of theorems.  If dimension
is itself an input, the data must instead be list-coded and the algorithm must be a single primitive-
recursive program over those codes.  Both are legitimate, but their resource and invariant proofs differ.
